\begin{lemma}
\label{lemma-cone-squared}
Let $R$ be a ring. Let $A_\bullet$ be a complex of $R$-modules.
Let $f, g \in R$. Let $C(f)_\bullet$ be the cone of
$f : A_\bullet \to A_\bullet$. Define similarly $C(g)_\bullet$ and
$C(fg)_\bullet$. Then $C(fg)_\bullet$ is homotopy equivalent to the
cone of a map
$$
C(f)_\bullet[1] \longrightarrow C(g)_\bullet
$$
\end{lemma}

\begin{proof}
We first prove this if $A_\bullet$ is the complex consisting of $R$ placed
in degree $0$. In this case the complex $C(f)_\bullet$ is the
complex
$$
\ldots \to 0 \to R \xrightarrow{f} R \to 0 \to \ldots
$$
with $R$ placed in (homological) degrees $1$ and $0$. The map
of complexes we use is
$$
\xymatrix{
0 \ar[r] \ar[d] &
0 \ar[r] \ar[d] &
R \ar[r]^f \ar[d]^1 &
R \ar[r] \ar[d] & 0 \ar[d] \\
0 \ar[r] & R \ar[r]^g & R \ar[r] & 0 \ar[r] & 0
}
$$
The cone of this is the chain complex consisting of $R^{\oplus 2}$ placed in
degrees $1$ and $0$ and differential (\ref{equation-differential-cone})
$$
\left(
\begin{matrix}
g & 1 \\
0 & -f
\end{matrix}
\right) :
R^{\oplus 2} \longrightarrow R^{\oplus 2}
$$
To see this chain complex is homotopic to
$C(fg)_\bullet$, i.e., to $R \xrightarrow{fg} R$,
consider the maps of complexes
$$
\xymatrix{
R \ar[d]_{(1, -g)} \ar[r]_{fg} & R \ar[d]^{(0, 1)} \\
R^{\oplus 2} \ar[r] & R^{\oplus 2}
}
\quad\quad
\xymatrix{
R^{\oplus 2} \ar[d]_{(1, 0)} \ar[r] & R^{\oplus 2} \ar[d]^{(f, 1)} \\
R \ar[r]^{fg} & R
}
$$
with obvious notation.
The composition of these two maps in one direction is the
identity on $C(fg)_\bullet$, but in the other direction
it isn't the identity. We omit writing out the required homotopy.

\medskip\noindent
To see the result holds in general, we use that we have a functor
$K_\bullet \mapsto \text{Tot}(A_\bullet \otimes_R K_\bullet)$
on the category of complexes which is compatible with homotopies
and cones. Then we write $C(f)_\bullet$ and $C(g)_\bullet$
as the total complex of the double complexes
$$
(R \xrightarrow{f} R) \otimes_R A_\bullet
\quad\text{and}\quad
(R \xrightarrow{g} R) \otimes_R A_\bullet
$$
and in this way we deduce the result from the special case discussed above.
Some details omitted.
\end{proof}

\begin{lemma}
\label{lemma-koszul-mult-abstract}
Let $R$ be a ring. Let $\varphi : E \to R$ be an $R$-module map.
Let $f, g \in R$. Set $E' = E \oplus R$ and define
$\varphi'_f, \varphi'_g, \varphi'_{fg} : E' \to R$
by $\varphi$ on $E$ and multiplication by $f, g, fg$ on $R$.
The complex $K_\bullet(\varphi'_{fg})$ is homotopy equivalent to the
cone of a map of complexes
$$
K_\bullet(\varphi'_f)[1]
\longrightarrow
K_\bullet(\varphi'_g).
$$
\end{lemma}

\begin{proof}
By
Lemma \ref{lemma-cone-koszul-abstract}
the complex $K_\bullet(\varphi'_f)$ is isomorphic to the cone of
multiplication by $f$ on $K_\bullet(\varphi)$ and similarly
for the other two cases. Hence the lemma follows from
Lemma \ref{lemma-cone-squared}.
\end{proof}

\begin{lemma}
\label{lemma-koszul-mult}
Let $R$ be a ring. Let $f_1, \ldots, f_{r - 1}$ be a sequence of elements
of $R$. Let $f, g \in R$. The complex
$K_\bullet(f_1, \ldots, f_{r - 1}, fg)$
is homotopy equivalent to the cone of a map of complexes
$$
K_\bullet(f_1, \ldots, f_{r - 1}, f)[1]
\longrightarrow
K_\bullet(f_1, \ldots, f_{r - 1}, g)
$$
\end{lemma}

\begin{proof}
Special case of
Lemma \ref{lemma-koszul-mult-abstract}.
\end{proof}
